\documentclass{nature}
\usepackage{amsmath}
\usepackage{amssymb}
\usepackage{graphicx}
\usepackage{booktabs}
\usepackage{biblatex}
\usepackage{hyperref}
\addbibresource{refs.bib}

\graphicspath{{./figures/}}

\title{Single-cell RNA sequencing reveals novel insights into cellular heterogeneity}

\author{First Author$^{1,2}$, Second Author$^{1,3}$, and Third Author$^{1,4}$}

\begin{document}

\begin{abst}
Single-cell RNA sequencing has revolutionized our understanding of cellular
heterogeneity in complex tissues. Here we present a comprehensive analysis of
cell types in mammalian brain tissue using scRNA-seq~\supercite{zheng2017}. We
identify 47 distinct cell populations and characterize their transcriptional
signatures. Our findings reveal previously unknown cell types and provide insights
into brain development and disease mechanisms.
\end{abst}

\maketitle

\section{INTRODUCTION}

The advent of single-cell RNA sequencing (scRNA-seq) has enabled researchers to
profile gene expression at unprecedented resolution~\supercite{tabula,regev2017}.
This technology has transformed our understanding of cellular diversity in
complex tissues.

\section{RESULTS}

We performed scRNA-seq on 50,000 cells from adult mouse brain tissue.

\begin{figure}[ht]
  \centering
  \includegraphics[width=0.9\textwidth]{fig-overview.pdf}
  \caption{Single-cell analysis overview. a, UMAP visualization of cell types.
    b, Marker gene expression heatmap. c, Cell type proportions.}
  \label{fig:overview}
\end{figure}

\begin{table}[ht]
\centering
\caption{Major cell type abundances in the dataset.}
\label{tab:celltypes}
\begin{tabular}{lcc}
  \toprule
  Cell Type & Count & Percentage \\
  \midrule
  excitatory neurons & 15,234 & 30.5\% \\
  inhibitory neurons & 8,567 & 17.1\% \\
  astrocytes & 6,789 & 13.6\% \\
  oligodendrocytes & 5,432 & 10.9\% \\
  microglia & 2,156 & 4.3\% \\
  other & 11,822 & 23.6\% \\
  \bottomrule
\end{tabular}
\end{table}

\section{DISCUSSION}

Our study reveals unprecedented cellular diversity in the mammalian brain.
The transcriptional profiles we characterized provide a foundation for
understanding brain function and disease.

\section{METHODS}

Cells were dissociated using papain and loaded onto the 10x Chromium platform.
Sequencing was performed on an Illumina NovaSeq 6000.

\printbibliography

\end{document}
